
\documentclass{article}
\usepackage{colortbl}
\usepackage{makecell}
\usepackage{multirow}
\usepackage{supertabular}

\begin{document}

\newcounter{utterance}

\centering \large Interaction Transcript for game `cladder', experiment `full\_v1.5\_default', episode 403 with Qwen3.6{-}27B.
\vspace{24pt}

{ \footnotesize  \setcounter{utterance}{1}
\setlength{\tabcolsep}{0pt}
\begin{supertabular}{c@{$\;$}|p{.15\linewidth}@{}p{.15\linewidth}p{.15\linewidth}p{.15\linewidth}p{.15\linewidth}p{.15\linewidth}}
    \# & \multicolumn{2}{c}{Player} && \multicolumn{2}{c}{Game Master} \\
    \hline

    \theutterance \stepcounter{utterance}  
    & & & \multicolumn{4}{p{0.6\linewidth}}{
        \cellcolor[rgb]{0.9,0.9,0.9}{
            \makecell[{{p{\linewidth}}}]{
                \texttt{\tiny{[P1$\langle$GM]}}
                \texttt{You are an expert in causal inference. The following question is not a typical commonsense query, but rather a meticulously designed question created by a professor specializing in causal inference, intended to assess the students' mastery of the course content.} \\
\texttt{Imagine a self{-}contained, hypothetical world with only the following conditions, and without any unmentioned factors or causal relationships: Vaccination status has a direct effect on vaccination reaction and getting smallpox. Getting smallpox has a direct effect on smallpox survival. Vaccination reaction has a direct effect on smallpox survival. For unvaccinated individuals, the probability of smallpox survival is 39\%. For vaccinated individuals, the probability of smallpox survival is 46\%. For vaccinated individuals, would it be more likely to see smallpox survival if the person had been vaccinated?} \\
\texttt{Start your answer with 'Yes' or 'No', followed by additional reasoning or evidence to support your explanation.} \\
            }
        }
    }
     \\ \\

    \theutterance \stepcounter{utterance}  
    & \multicolumn{4}{p{0.6\linewidth}}{
        \cellcolor[rgb]{0.561,0.737,0.561}{
            \makecell[{{p{\linewidth}}}]{
                \texttt{\tiny{[P1$\rangle$GM]}}
                \texttt{No} \\
\\ 
\texttt{**Reasoning:**} \\
\\ 
\texttt{To answer this question, we must distinguish between the **observed association** (the difference in survival rates between vaccinated and unvaccinated groups) and the **causal effect** of vaccination on survival for a specific individual who is already vaccinated.} \\
\\ 
\texttt{1.  **Analyze the Causal Structure:**} \\
\texttt{    *   \$V\$ (Vaccination) \$\textbackslash{}rightarrow\$ \$R\$ (Reaction)} \\
\texttt{    *   \$V\$ (Vaccination) \$\textbackslash{}rightarrow\$ \$S\$ (Smallpox Infection)} \\
\texttt{    *   \$S\$ (Smallpox Infection) \$\textbackslash{}rightarrow\$ \$Y\$ (Survival)} \\
\texttt{    *   \$R\$ (Reaction) \$\textbackslash{}rightarrow\$ \$Y\$ (Survival)} \\
\\ 
\texttt{    The question asks: "For vaccinated individuals, would it be more likely to see smallpox survival if the person had been vaccinated?"} \\
\\ 
\texttt{2.  **Interpret the Counterfactual:**} \\
\texttt{    The phrase "if the person had been vaccinated" implies a counterfactual comparison. However, the subject is already defined as a "vaccinated individual."} \\
\texttt{    *   Let \$Y(1)\$ be the survival outcome if vaccinated.} \\
\texttt{    *   Let \$Y(0)\$ be the survival outcome if unvaccinated.} \\
\texttt{    *   The question asks about the causal effect for a specific unit \$i\$ who is vaccinated (\$V\_i=1\$). It} \\
            }
        }
    }
    & & \\ \\

    \theutterance \stepcounter{utterance}  
    & & & \multicolumn{2}{p{0.3\linewidth}}{
        \cellcolor[rgb]{0.9,0.9,0.9}{
            \makecell[{{p{\linewidth}}}]{
                \texttt{\tiny{[GM$|$GM]}}
                \texttt{No} \\
            }
        }
    }
    & & \\ \\

    \theutterance \stepcounter{utterance}  
    & & & \multicolumn{2}{p{0.3\linewidth}}{
        \cellcolor[rgb]{0.9,0.9,0.9}{
            \makecell[{{p{\linewidth}}}]{
                \texttt{\tiny{[GM$|$GM]}}
                \texttt{no} \\
            }
        }
    }
    & & \\ \\

    \theutterance \stepcounter{utterance}  
    & & & \multicolumn{2}{p{0.3\linewidth}}{
        \cellcolor[rgb]{0.9,0.9,0.9}{
            \makecell[{{p{\linewidth}}}]{
                \texttt{\tiny{[GM$|$GM]}}
                \texttt{game\_result = WIN} \\
            }
        }
    }
    & & \\ \\

\end{supertabular}
}

\end{document}
